\thispagestyle{plain}
\chapter*{Resumen}

El diseño y descubrimiento de fármacos está condicionado por el tamaño del espacio químico estimado en más de $10^{60}$ moléculas posibles, esta escala resulta inviable de explorar mediante métodos experimentales tradicionales o de cribado virtual exhaustivo. Además, el explorar y diseñar compuestos con características de fármaco es un problema esencialmente multiobjetivo, ya que exige optimizar simultáneamente propiedades fisicoquímicas y farmacológicas frecuentemente en conflicto.

Frente a estas limitaciones, los modelos generativos permiten representar dicho espacio molecular en un espacio latente continuo, mientras que la aplicación de algoritmos de optimización multiobjetivo para guiar la búsqueda de nuevos compuestos en este espacio permite dirigir esta generación hacia candidatos con propiedades deseables.

Esta tesina de grado presenta la implementación y evaluación de un enfoque híbrido que combina autoencoders variacionales junto con algoritmos metaheurísticos multiobjetivo para optimizar la generación de nuevos compuestos explorando este espacio latente. Como caso de aplicación, este trabajo se orienta a la obtención de moléculas novedosas y viables para el tratamiento de la resistencia a fármacos en el cáncer colorrectal, relacionada con proteínas transportadoras de la familia ABC.


\clearpage

\thispagestyle{plain}
\chapter*{Abstract}

Drug design and discovery is conditioned by the size of the chemical space, estimated at more than $10^{60}$ possible molecules, making exploration through traditional experimental methods or exhaustive virtual screening infeasible. Furthermore, exploring and designing drug-like compounds is essentially a multi-objective problem, as it requires simultaneously optimizing physicochemical and pharmacological properties that are frequently in conflict.

To address these limitations, generative models allow representing this molecular space in a continuous latent space, while the application of multi-objective optimization algorithms to guide the search for new compounds within this space allows steering this generation toward candidates with desirable properties.

This bachelor's thesis presents the implementation and evaluation of a hybrid approach combining variational autoencoders with multi-objective metaheuristic algorithms to optimize the generation of new compounds by exploring this latent space. As an application case, this work focuses on obtaining novel and viable molecules for the treatment of drug resistance in colorectal cancer, associated with ABC family transporter proteins.

